\section{Stronger-engine characterization of simpler heads}
\label{app:baseline-engine}
After observing the additional teacher-agreement comparison, we froze another
post-hoc engine assessment. It includes every paired backbone and full-transport
policy plus all nine newly trained simpler heads. The original 128-position
subsets remain fixed; no policy or position is selected from favorable new
outcomes. Every unrestricted and unique chosen-move call starts with a cleared
hash and requests 20,000 nodes under the same single-threaded Stockfish setup.

\begin{table}[ht]
\centering\small
\caption{Post-hoc engine assessment of all fifteen simpler-head policies on the same fixed 128 positions per panel. All calls are fresh and request 20,000 nodes. Both loss measures are lower-is-better; raw centipawn loss includes an artificial mate-score encoding.}
\begin{tabular}{lrrrr}
\toprule
Head & Bounded, ordinary & Bounded, shifted & CP, ordinary & CP, shifted \\
\midrule
Frozen direct & 0.09878 & 0.08163 & 642.6 & 428.6 \\
Full transport & 0.06535 & 0.07278 & 619.7 & 451.8 \\
Generic MLP & 0.09502 & 0.08370 & 650.9 & 437.2 \\
NCN adaptation & 0.08712 & 0.08185 & 601.6 & 460.4 \\
Typed overlap & 0.08773 & 0.08364 & 600.7 & 423.3 \\
\bottomrule
\end{tabular}
\end{table}


\BaselineEngineResults{}

\paragraph{Published path-baseline adaptation.}
A separate three-fit study uses a 16,416-parameter NBFNet-style head: a
120-to-16 action query, three width-16 DistMult/PNA layers with source boundary
reinjection, residual shortcuts and layer normalization, followed by a target
lookup and width-32 readout. The zero-initialized scalar correction preserves
the same original backbones. It uses the same four typed root attack channels,
32,768 roots, six epochs and 1,536 updates per seed; all fits precede evaluation.
Selected pinned public unfused methods pass value/gradient checks through a
scatter compatibility layer. This is a documented chess adaptation, not an
official NBFNet benchmark reproduction. Unlike transport, it accesses frozen
node features through the endpoint/global query rather than pooling features
of every visited node. Similar stored parameters do not equalize compute.

\begin{table}[ht]
\centering\small
\caption{Separately frozen path-baseline comparison. Agreement is percent; NLL is lower-is-better. All three paired seeds retained; original references copied with authenticated provenance.\label{tab:path-baseline}}
\begin{tabular}{lrrrr}
\toprule
Head & Agree., ordinary & Agree., shifted & NLL, ordinary & NLL, shifted \\
\midrule
Frozen direct (copied) & 31.75 & 26.81 & 2.435 & 2.635 \\
Transport (copied) & 33.02 & 27.99 & 2.358 & 2.575 \\
NBFNet-style path head & 34.54 & 28.91 & 2.298 & 2.514 \\
\bottomrule
\end{tabular}
\end{table}

The original transport needed at least one agreement point over this comparator
on both panels, with no paired-seed deficit worse than half a point. Both checks
fail. The audit verifies all 4,608 update-index records and replays 12,288 new
baseline predictions plus 24,576 copied references. The original failed
transport criteria remain unchanged; there is no new engine or gameplay result.
\PathBaselineLatency{}

\section{Candidate-specific attack operators}
Root-graph transport omits changes to attack lines caused by a proposed move.
An untrained follow-up computes four child operators in the fixed root-player
frame, using native transitions for all legal candidates. Write
$D_{a,r}=P_{\mathrm{child}(a),r}-P_r$, including all changed normalization
coefficients and empty-row self-loops. Sparse indexed accumulation implements
$pP_r+pD_{a,r}$. This is the linear identity for applying the child operator,
not a new theorem or learned transition model. Both paths first reconstruct
native child graphs; the prototype is not an incremental attack generator.
Nine tests cover color frames, castling, en passant, promotions, dense operator
reconstruction and three-step output/gradient equivalence.

The frozen screen uses point masses at the move endpoints and seeded random
gates, node features and readout weights. It has no trained query controller or
chess-quality evaluation. Complete untrained pipeline timing includes board
parsing, legal enumeration, graph construction, normalization/delta packing,
input preparation, three propagation steps, readout and UCI selection.
It uses sixteen roots, three paired timing repeats and a shared host.
\CounterfactualResults{}
Any later repair must preserve this failed screen. A quality study would first
need to establish a benefit from the full child-operator representation against
root-only and candidate-assignment controls.

\paragraph{Shared-weight graph contrast.}
A separate dense prototype scores the change in a candidate-conditioned graph
representation. With frozen root features $H_s$, action features $q_s(a)$ and
the same learned head $f_\theta$ in both branches, its correction is
\begin{equation}
 c(s,a)=f_\theta(P_{s,a};H_s,q_s(a))-f_\theta(P_s;H_s,q_s(a)).
\end{equation}
The child graph comes from the exact native move in the root-player frame;
node features remain those of the root. Shared controllers may diverge after
receiving different pooled contexts, so this is not the fixed-gate linear
operator identity above. A zero-change mask enforces exact cancellation in
both forward and backward execution when the operators agree. The root term
is action-dependent, so subtracting it can affect the legal-move ranking.

The four implemented arms are this difference, child only, root only, and a
label-independent permutation of child-to-candidate assignments. Each stores
16,744 parameters and starts with zero residual output. Eighteen tests include
special moves, direction, parameter identity, separate-branch equivalence and
menu alignment. The first screen stopped on a tuple/list validation mismatch;
its failed receipt is retained. A separately frozen repair changes that
validation, retaining all positions, seeds and thresholds. \GraphContrastResults{}
\GraphContrastSharedResults{}

\paragraph{Completed scalar-contrast quality study.}
\GraphQualityResult{}
\begin{table}[ht]
\centering\small
\caption{Completed scalar-contrast study on ordinary (O) and shifted (S) panels. Agreement is percent; NLL is lower-is-better. All three seeds retained. Corrupted contrast reuses fitted contrast weights; it is a diagnostic, not a new fit.\label{tab:graph-contrast}}
\begin{tabular}{lrrrr}
\toprule
Head & Agree. O & Agree. S & NLL O & NLL S \\
\midrule
Frozen base & 31.75 & 26.81 & 2.435 & 2.635 \\
Root only & 33.02 & 27.99 & 2.358 & 2.575 \\
Child only & 33.25 & 28.14 & 2.347 & 2.573 \\
Scalar contrast & 33.41 & 27.95 & 2.364 & 2.571 \\
Trained shuffled contrast & 32.06 & 26.33 & 2.432 & 2.637 \\
Test-time shuffled contrast & 31.33 & 26.68 & 2.458 & 2.647 \\
\bottomrule
\end{tabular}
\end{table}


\AttackEditInputResult{}
\NormalizationInputResult{}

Paired chess representations already use shared weights \citep{david2017}.
More directly, WLDN ranks candidate reaction outcomes using differences of graph
representations and adds its score to an earlier proposal score
\citep{coley2019}. Shared graph encoders with subtractive atom representations
also predict activation energies \citep{grambow2020}, while condensed reaction
graphs can encode changes directly in node and edge features \citep{heid2022}.
Our action-conditioned transport and scalar subtraction differ from their
nodewise difference encoders; this distinction alone is not a contribution.
The ingredient controls do not replace the direct method comparisons below.
Perturbed-graph explanation is another established use \citep{lucic2022}.

\paragraph{Direct prior-method adapters.}
We implement two additional baselines to separate scalar-score differences,
nodewise differences and changes encoded directly on the input graph.
The WLDN-style adapter \citep{coley2019} applies a shared width-32 encoder with
three tied WL updates to root and child graphs, subtracts the node states, and
applies one separate WL update on the child graph before summing over squares.
It concatenates the result with the original 120 action features and uses a
59-unit readout, for 16,638 head parameters. Edge features and biases in the
difference encoder mean identical branches need not have zero learned scores.
We retain that property rather than imposing scalar-contrast cancellation.

The condensed-graph adapter \citep{heid2022} uses the union of root and child
attack edges. Each message direction carries four root direction/color flags
and four signed child-minus-root flags. A directed MPNN excludes immediate
reverse-edge messages and adds the original edge projection at each of two
tied recurrent updates. Width is 32, pooling averages all 64 square states,
and an 80-unit readout receives the same action features, for 16,704 parameters.
All node features remain those of the root in both adapters; the resulting
zero node differences are omitted from the condensed input. These are explicit
chess adaptations, not reproductions of molecular features or chemical tasks.

Selected pinned public functions pass output and gradient checks: WLDN through
a Torch compatibility layer, and the Chemprop directed encoder through its
unmodified Torch class with supplied graph tensors. The WLDN source and port
retain GPL identification; the selected Chemprop source and adapter retain MIT
identification. Neither check reproduces the full published training pipeline.
\DifferenceBaselineEngineering{}
A separately frozen three-seed WLDN study uses the same 32,768 roots, six epochs
and minibatch order as scalar contrast. \WLDNQualityResult{}
A separate condensed-graph study completes the same three seeds and 4,608-update
budget, with all 24,576 new/copied predictions replayed. Its protocol was frozen
before inspecting final WLDN or graph-contrast quality outcomes, with NBFNet's
advantage over original transport already known.

\paragraph{Combined established-method comparison.}
\label{app:graph-baselines}
The first separate comparison retains all eight original ingredient checks and
adds two WLDN checks. A further frozen comparison adds two checks each for CGR
and NBFNet. Every new check requires at least one point of mean contrast
advantage and no paired-seed deficit beyond half a point. All fourteen checks
are required; a later comparison cannot rescue an earlier failure. WLDN and
NBFNet outcomes were known when the combined analysis was frozen, while scalar
contrast and CGR were still training. This is adaptive development, not a
preregistered research program.
\begin{table}[ht]
\centering\small
\caption{Identity-joined graph-baseline comparison on the same ordinary (O) and shifted (S) development panels. Agreement is percent; NLL is lower-is-better. Milliseconds are median complete native decisions in one balanced six-method process. Three paired seeds; similar parameter counts do not match inference compute.\label{tab:graph-baselines}}
\begin{tabular}{lrrrrr}
\toprule
Head & Agree. O & Agree. S & NLL O & NLL S & ms \\
\midrule
Frozen base & 31.75 & 26.81 & 2.435 & 2.635 & 0.452 \\
Root transport & 33.02 & 27.99 & 2.358 & 2.575 & 0.889 \\
NBFNet-style & 34.54 & 28.91 & 2.298 & 2.514 & 4.248 \\
Scalar contrast & 33.41 & 27.95 & 2.364 & 2.571 & 3.734 \\
Nodewise WLDN & 37.01 & 31.43 & 2.161 & 2.355 & 4.876 \\
Condensed graph & 34.08 & 28.74 & 2.333 & 2.542 & 3.962 \\
\bottomrule
\end{tabular}
\end{table}

\FloatBarrier
\GraphBaselineResult{}

\FloatBarrier
\section{Fixed-weight WLDN pathway diagnostic}
\label{app:wldn-pathways}
After inspecting the completed baseline comparisons, we froze a diagnostic of
the strongest adapted head. It retains all three trained WLDN checkpoints and
both complete 2,048-position development panels. No weights, labels or legal
menus change. Five interventions separate the encoded child-minus-root node
differences from the graph inputs to the subsequent difference layer.

\begin{table}[ht]
\centering\small
\caption{Fixed-weight WLDN pathway interventions on ordinary (O) and shifted (S) development panels. All three seeds retained. Agreement is percent; NLL is lower-is-better. These are interventions on trained heads, not separately trained architectures.\label{tab:wldn-interventions}}
\begin{tabular}{lrrrr}
\toprule
Intervention & Agree. O & Agree. S & NLL O & NLL S \\
\midrule
Frozen base reference & 31.75 & 26.81 & 2.435 & 2.635 \\
Native WLDN & 37.01 & 31.43 & 2.161 & 2.355 \\
Zero node difference & 31.56 & 26.25 & 2.436 & 2.647 \\
Root graph at difference layer & 36.30 & 30.00 & 2.204 & 2.411 \\
Zero difference-layer edge flags & 35.07 & 29.49 & 2.243 & 2.432 \\
Zero pooled graph channels & 31.25 & 26.68 & 2.432 & 2.623 \\
Shuffle node differences & 26.38 & 22.28 & 2.826 & 3.056 \\
\bottomrule
\end{tabular}
\end{table}


Zeroing node differences retains child adjacency, edge flags and learned
biases. Replacing the difference-layer graph with the root graph retains the
native encoded differences. Zeroing its four edge flags retains child
adjacency, so this intervention does not remove all structural information.
Zeroing the pooled graph channels retains the original 120 action features
and trained readout. Finally, the frozen candidate permutation shuffles node
differences within each legal menu while keeping the actual child's graph and
action features. The permutation can contain fixed points.

\WLDNInterventionResult{}
The audit also verifies unchanged head states, removal of temporary hooks,
native root observations, legal menus and labels. It authenticates the prior
complete native child-graph audit rather than reconstructing that cache again.
Bootstrap resampling uses 2,000 draws of source games after averaging paired
correctness across the three seeds; the same draws serve all interventions
within a panel.

The trained model depends on candidate-aligned node changes and on the
subsequent graph context. These effects are not additive contributions:
interventions can create out-of-distribution hidden states, and the difference
layer combines node and edge information nonlinearly. In particular, the
zero-difference condition is not a separately trained edge-only architecture.
These results motivate retaining nodewise differences in the next controlled
architecture comparison; they do not establish a new mechanism or predict the
performance of retrained variants.

\FloatBarrier
\section{Union difference-layer hypothesis and encoding limit}
\label{app:union-basis}
An implemented extension retains the shared WLDN encoder \citep{jin2017} and nodewise
difference but changes the graph used by the final difference layer. Four arms
use child connectivity, root/child union connectivity, a union with explicit
retained/added/removed relation flags, or the same union with corrupted edit
labels. Each head stores 16,740 parameters with a width-32 encoder and a
58-unit action readout. The child/union controls leave eight of twelve edge
input channels unused, so stored-parameter matching does not equalize active
capacity. \UnionDifferenceEngineering{}

The negative control rotates edit classes over ordered square pairs within
each candidate and direction/color channel, preserving union flags and class
counts. It can break reverse-direction consistency and leave homogeneous
groups unchanged. A globally fixed added/removed swap would instead be a
learnable input renaming. The separate quality study trains all twelve final
models before evaluation, retaining the original WLDN and backbone as
authenticated references. Its unchanged criterion requires an edits advantage
over every trained comparator on both panels.

\UnionEditControlInput{}

\paragraph{Matched training result.}
\UnionDifferenceQuality{}
\begin{table}[t]
\centering\small
\caption{Union difference-layer development comparison. All three paired seeds retained. Agreement is percent; NLL is lower-is-better. Time is the median paired complete-native decision ratio to WLDN on sixteen fixed ordinary roots, with six rotating repeats. The test-time corruption reuses edits weights and has no separate timing. \label{tab:union-difference}}
\begin{tabular}{lrrrrr}
\toprule
Policy & Agree., ord. & Agree., shift & NLL, ord. & NLL, shift & Time / WLDN \\
\midrule
Frozen direct & 31.75 & 26.81 & 2.435 & 2.635 & 0.112 \\
Original WLDN & 37.01 & 31.43 & 2.161 & 2.355 & 1.000 \\
Matched child & 36.88 & 31.15 & 2.172 & 2.363 & 1.414 \\
Union & 36.70 & 31.72 & 2.163 & 2.365 & 1.437 \\
Edit labels & 37.22 & 31.92 & 2.153 & 2.356 & 1.510 \\
Rotated labels & 36.64 & 31.53 & 2.161 & 2.358 & 1.626 \\
Edits, corrupted at test & 36.36 & 31.36 & 2.167 & 2.366 & -- \\
\bottomrule
\end{tabular}
\end{table}

\UnionDifferenceCost{}

\paragraph{Reverse-consistent control.}
\UnionConsistentEngineering{}

\paragraph{Equivalent feature coordinates.}
For one binary relation, let $r,c,u=r\lor c$ denote root, child and union
indicators, and let $t,a,\ell$ denote retained, added and removed indicators.
Then
\begin{equation}
 t=r+c-u,\qquad a=u-r,\qquad \ell=u-c;
 \qquad u=t+a+\ell,\quad r=t+\ell,\quad c=t+a.
\end{equation}
Thus these three-channel encodings are invertibly related for each of the four
direction/color relations. An edge projection with weights $W_t,W_a,W_\ell$
has exactly the same preactivation in the presence basis when
\begin{equation}
 W_u=-W_t+W_a+W_\ell,\qquad W_r=W_t-W_a,\qquad W_c=W_t-W_\ell.
\end{equation}
All node columns, biases, connectivity and subsequent layers remain unchanged,
so the complete scoring function is preserved in real arithmetic. Seven tests
verify the Boolean truth table, inverse transforms, full-head scores and
chain-rule gradients, including nonzero projections and hook cleanup.
The float64 score/gradient tolerances are $10^{-12}/10^{-10}$.

This is an elementary reparameterization, not a new theorem or an expressivity
advantage. It does not equate child-only with union connectivity. The transform
is nonorthogonal, so isotropic random initialization, coordinatewise Adam and
Euclidean gradient clipping need not produce equivalent training trajectories.
Any basis-dependent training benefit would require a separate controlled
comparison. Established condensed reaction graphs already encode both sides
of a transformation \citep{heid2022}; naming edit classes is insufficient for
an algorithmic novelty claim.
BonDNet also forms learned atom/bond differences and represents a broken bond
by a negative reactant feature before set-to-set pooling \citep{wen2021}.
Its readout differs from the final union-graph WL step here; that distinction
alone does not establish novelty.

\section{Selected-panel overlap sensitivity}
A post-hoc input audit connects source games sharing a selected root or a legal
one-move successor, treating the color/rank mirror as equivalent and ignoring
move counters while retaining literal en-passant state. The ordinary panel
retains 39 components for 39 games, and the shifted panel retains 31 for 31;
there are no cross-game links under this exposure definition. Largest groups
contain 61 and 92 positions. Thus this particular overlap grouping does not
change source-game bootstrap units. It does not certify whole-rollout
separation, general statistical independence, seed-population uncertainty or
protection against adaptive use of the development panels. No new positions,
model predictions or engine calls are used for this audit.
